\section{Proposed Framework}
\label{sec:framework}

This section develops the verification and recovery model in full.
Fig.~\ref{fig:pipeline} shows how the pieces connect; notation is summarised
in Table~\ref{tab:notation}.

\begin{figure}[!t]
\centering
\begin{tikzpicture}[
  font=\footnotesize,
  box/.style={draw=black!70, rounded corners=2pt, align=center, fill=gray!8,
              minimum width=4.6cm, inner sep=4pt},
  ev/.style={box, fill=white, minimum width=3.6cm},
  st/.style={draw=black!70, rounded corners=2pt, minimum width=2.4cm,
             inner sep=4pt, fill=white},
  iso/.style={box, fill=orange!8, draw=orange!60!black, minimum width=3.9cm},
  note/.style={align=left, font=\scriptsize, text width=3.6cm},
  arr/.style={-{Stealth[length=2mm]}, black!70},
  node distance=6mm]
\node[ev] (beh) {\textbf{Behavioural evidence}\\forwarding ratio, peer-relative $z_B$};
\node[ev, right=8mm of beh] (loc) {\textbf{Localisation evidence}\\range residual, calibrated $z_L$};
\node[box, below=12mm of $(beh)!0.5!(loc)$] (s1) {\textbf{Stage 1: weighted Stouffer fusion}\\
  $Z_{ij}=(\tilde\omega_B A_B+\tilde\omega_L A_L)/\sqrt{\tilde\omega_B^2+\tilde\omega_L^2}$};
\node[box, below=of s1] (td) {Direct trust $T^D_{ij}=\exp(-Z_{ij}^2/2)$, per observer};
\node[box, below=of td] (s2) {\textbf{Stage 2: aggregation}\\
  $T_{ij}=(w_D T^D+w_I T^I+w_H T^H)/\sum w$};
\node[box, below=of s2] (dec) {Median over observers, agreement, persistence $M_i \ge K$};
\node[st, below=9mm of dec] (sus) {Suspicious};
\node[st, left=8mm of sus] (nor) {Normal};
\node[st, right=8mm of sus] (con) {Confirmed};
\node[iso, below=8mm of con] (iset) {Isolation set $\mathcal{B}_{iso}(t)$\\blacklist queueing discipline};
\node[iso, below=5mm of iset] (rec) {Route repair (AODV)\\$P \cap \mathcal{B}_{iso} = \emptyset$};
\draw[arr] (beh.south) |- ([yshift=5mm]s1.north) -- (s1.north);
\draw[black!70] (loc.south) |- ([yshift=5mm]s1.north);
\draw[arr] (s1) -- (td);
\draw[arr] (td) -- (s2);
\draw[arr] (s2) -- (dec);
\draw[arr] (dec.south) -- ++(0,-4mm) -| (nor.north);
\draw[arr] (dec.south) -- (sus.north);
\draw[arr] (dec.south) -- ++(0,-4mm) -| (con.north);
\draw[arr, orange!60!black] (con) -- (iset);
\draw[arr, orange!60!black] (iset) -- (rec);
\draw[arr, dashed, orange!60!black] (rec.west) -| node[pos=0.3, above, font=\scriptsize]{release if revised (E, E$'$)} (nor.south);
\node[note, right=3mm of s1] {Unit variance under the null for any weights, so one threshold serves all configurations (Sec.~\ref{sec:calib}).};
\node[note, right=3mm of s2] {$T^I$: trimmed peer mean, one recommender capped at $\kappa=0.3$. $T^H$: asymmetric EWMA, $\lambda_\downarrow=0.5$, $\lambda_\uparrow=0.9$.};
\node[note, right=3mm of iset] {Policies D, D$_w$, D$_1$ isolate before confirmation and never release; E and E$'$ isolate only from Confirmed and can release.};
\end{tikzpicture}
\caption{Verification pipeline. Behavioural and localisation evidence are
standardised, oriented and fused in Stage~1 to give a per-observer direct
trust value; Stage~2 aggregates that value across observers and over time.
The three-state verdict separates suspicion from confirmation, and the
isolation policies differ only in which evidence authorises removal and
whether removal can be reversed.}
\label{fig:pipeline}
\end{figure}

\subsection{Behavioural evidence}

For beacon $i$ in window $w$, let $P_i^{fwd}$ and $P_i^{recv}$ denote packets
forwarded and received for forwarding. The behavioural observable is the
forwarding ratio
\begin{equation}
\label{eq:fr}
FR_i(t) = \frac{P_i^{fwd}(t)}{P_i^{recv}(t) + \epsilon},
\end{equation}
with $\epsilon$ guarding against division by zero.

We use a single behavioural observable rather than a composite score, because
composites in this literature are often degenerate. A weighted sum of the
form $w_1 FR_i + w_2(1-DR_i) + \dots$, where $DR_i = 1 - FR_i$ is the drop
ratio, collapses to $(w_1 + w_2) FR_i$: what reads as two independent pieces
of evidence is one observable counted twice, with an inflated effective
weight. Observables genuinely independent of $FR_i$, such as routing-response
consistency or response delay, would each require their own null calibration
and are outside the scope of this study.

The observable is standardised against contemporaneous peers rather than a
static calibrated baseline. Benign congestion depresses forwarding across all
nodes in a neighbourhood simultaneously; measured against a fixed baseline
this registers as individual misbehaviour, whereas measured against peers
observed in the same window a network-wide depression cancels. If
$\mathcal{P}_i$ is the peer set of $i$ with at least $n_{min}$ members,
\begin{equation}
\label{eq:zb}
z_B(w) = \frac{FR_i(w) - \mu(\mathcal{P}_i, w)}{\sigma(\mathcal{P}_i, w) + \epsilon},
\end{equation}
falling back to the static calibrated null where the peer set is too small to
estimate. Equation~\eqref{eq:zb} is the quantity that enters fusion.

\subsection{Localisation evidence}

\subsubsection{Range residual}

The detector is not given the beacon's true position: a node in a MANET has no
such oracle, and a model that consumes one is not implementable. We therefore
score the internal inconsistency of what the beacon broadcasts, which requires
only quantities an observer can measure.

For observer $j$ receiving a beacon from $i$, let $\hat d_{ij}$ be the distance
implied by measured RSSI under the calibrated path-loss model, and
$\tilde d_{ij}$ the distance implied by the coordinates the beacon advertises
together with $j$'s own position estimate. The range residual is
\begin{equation}
\label{eq:resid}
r_{ij}(w) = \bigl| \hat d_{ij}(w) - \tilde d_{ij}(w) \bigr|,
\end{equation}
and the standardised localisation evidence is $r_{ij}$ scored against the
residual null distribution calibrated on clean-network runs,
\begin{equation}
\label{eq:zl}
z_L(w) = \frac{r_{ij}(w) - \mu_r}{\sigma_r + \epsilon}.
\end{equation}
A beacon displaced by $\delta$ inflates $\tilde d_{ij}$ without changing
$\hat d_{ij}$, so $r_{ij}$ grows with $\delta$; an honest beacon leaves $r_{ij}$
at the measurement noise floor. True positions appear in our logs only for
labelling detections as correct or incorrect after the fact.

Two consequences follow. First, the residual is bounded below by RSSI ranging
error, so the attack is undetectable for offsets comparable to that error; the
detection floor is a property of the ranging channel and not of the fusion
rule. Second, an attacker who could also manipulate the RSSI observed by its
neighbours could suppress the residual. Our threat model does not include
physical-layer power manipulation, and we leave it as an open gap.

\subsection{Stage one: evidence fusion}

\subsubsection{Orienting the two evidence streams}

The two standardised statistics point in opposite directions with respect to
misbehaviour. A malicious beacon depresses its forwarding ratio, so $z_B$
becomes large and negative; it inflates its range residual, so $z_L$ becomes
large and positive. Summing the two as they stand would allow a node
misbehaving on both axes to cancel its own evidence. We therefore define a
common one-sided anomaly scale on which larger means more suspicious,
\begin{equation}
\label{eq:orient}
A_B(w) = -z_B(w), \qquad A_L(w) = +z_L(w),
\end{equation}
and all subsequent fusion operates on $A_B$ and $A_L$. The sign convention is
trivial once stated, but Section~\ref{sec:orientation} shows that getting it
wrong produces a detector that fails specifically on hybrid attacks while
continuing to work on single-axis ones.

\subsubsection{Confidence-weighted fusion}

Each branch is weighted by a fixed branch weight and by a confidence factor
saturating in sample count,
\begin{equation}
\label{eq:conf}
c(n) = \frac{n}{n + n_{1/2}},
\end{equation}
with $n_{1/2} = 5$ for localisation and $n_{1/2} = 45$ for behaviour. A single
beacon reception carries positional information directly, whereas forwarding
behaviour is observable only through an accumulation of transit packets. This
asymmetry reappears in the measured confirmation latencies of
Section~\ref{sec:results}. Writing the effective weights as
\begin{equation}
\label{eq:effw}
\tilde\omega_B(w) = \omega_B\, c(n_B(w)), \qquad
\tilde\omega_L(w) = \omega_L\, c(n_L(w)),
\end{equation}
where $n_B$ and $n_L$ are the observation counts available to each branch in
window $w$, the fused statistic is
\begin{equation}
\label{eq:stouffer}
Z_{ij}(w) = \frac{\tilde\omega_B A_B(w) + \tilde\omega_L A_L(w)}
                 {\sqrt{\tilde\omega_B^2 + \tilde\omega_L^2}}.
\end{equation}
A branch with no observations in a window has $c = 0$ and drops out, leaving
the remaining branch unchanged; this is why the confidence factor appears
inside the normalisation.

Direct trust is the fused statistic mapped through the Gaussian kernel
\begin{equation}
\label{eq:phi}
T^D_{ij}(w) = \phi\bigl(Z_{ij}(w)\bigr) = \exp\bigl(-\tfrac12 Z_{ij}(w)^2\bigr),
\end{equation}
so that $T^D = 1$ for a beacon indistinguishable from the calibrated null and
$T^D \to 0$ as the fused statistic moves into either tail. Thresholds on $T^D$
therefore correspond to thresholds on $|Z|$: the values used here,
$\tau_M = 0.40$ and $\tau_N = 0.70$, correspond to $|Z| = 1.354$ and
$|Z| = 0.845$, two-sided null tail probabilities of 0.176 and 0.398.

Two properties of~\eqref{eq:phi} matter. It is even in $Z$, so a deviation in
either direction reduces trust equally; this is deliberate, since a beacon
forwarding implausibly more traffic than its peers is as anomalous as one
forwarding less. And because $\phi$ is applied after fusion rather than before
it, the two-sidedness does not rescue a sign error: it is the fused $Z$ that is
folded, not the individual branches.

\subsubsection{Calibrating the fused statistic}
\label{sec:calib}

The combination rule matters less than how its output is thresholded. Let
$s_B, s_L \in [0,1]$ be evidence scores and $C = \omega_B s_B + \omega_L s_L$
with $\omega_B + \omega_L = 1$. Then
\begin{equation}
\label{eq:convex}
\min(s_B, s_L) \le C \le \max(s_B, s_L),
\end{equation}
so if both scores lie on the benign side of a threshold~$\tau$, so does $C$. A
detector that applies the single-branch threshold to a convex combination
therefore cannot confirm an attacker that both branches miss. Tables built
this way will still show the fused row ahead of the weaker branch, because a
mean of a strong score and a weak one exceeds the weak one; the apparent
advantage comes from the comparison and not from the mechanism.

The bound is a consequence of reusing $\tau$, not of averaging as such. For two
independent standardised branches the arithmetic mean
$\bar z = (z_B + z_L)/2$ has standard deviation $1/\sqrt2$ under the null, so a
threshold recalibrated to the same false-positive rate is $\tau/\sqrt2$, and
the rule $\bar z > \tau/\sqrt2$ is identical to the equal-weight Stouffer rule
$(z_B + z_L)/\sqrt2 > \tau$. What the fusion hypothesis requires is a combined
statistic calibrated against its own null distribution. Weighted Stouffer
combination~\eqref{eq:stouffer} is unit variance under the null whatever the
weights, so a single threshold serves every configuration and the hypothesis
becomes falsifiable: two branches each at $z = 1.5$ combine, with equal
weights, to $Z = 2.12$.

This argument disqualifies uncalibrated combination; it does not make
Stouffer's method the best calibrated rule, and we do not claim that it is.
Fisher's method, likelihood-ratio and Bayesian fusion, Dempster--Shafer
combination~\cite{singh2020self}, and learned rules can all be calibrated in
the same sense. We chose Stouffer's method because it accepts continuous
evidence without a likelihood model for either branch, admits per-branch
weights directly~\cite{whitlock2005weighted}, and degrades gracefully when one
branch has no observations in a window, which is frequent under mobility.
Section~\ref{sec:sweep} compares it against a disjunctive rule; we did not
compare it with the other calibrated rules.

Equation~\eqref{eq:stouffer} assumes $z_B$ and $z_L$ are independent under the
null. In our attack scenarios they are close to independent by construction,
because A1 and A2 disturb disjoint observables. Under a correlated failure,
such as congestion coinciding with a mobility-induced ranging error, the fused
statistic would be optimistic; correcting for it requires an estimated
correlation term.

\subsection{Stage two: trust aggregation}

\subsubsection{Direct trust}

Direct trust is the Stage~1 output $T^D_{ij}(w)$ of~\eqref{eq:phi}, evaluated
per observer per window. It is not separately smoothed; temporal weighting enters the model once, in
$T^H$ below.

\subsubsection{Indirect trust}

Let $R_{ij}(t) \in [0,1]$ be neighbour $j$'s recommendation regarding beacon
$i$ and $T_j(t)$ the reliability of that recommender:
\begin{equation}
\label{eq:ti}
T^I_i(t) = \frac{\sum_{j \in N_i} T_j(t) R_{ij}(t)}{\sum_{j \in N_i} T_j(t) + \epsilon}.
\end{equation}
Equation~\eqref{eq:ti} is the general reliability-weighted form. The evaluated
implementation uses a specialisation in which $R_{ij}$ is the direct trust
observer $j$ holds in target $i$, and the reliability weights are replaced by a
robust aggregation rule: a trimmed mean over the direct trust of observers
other than $j$, with $j$'s own assessment excluded to avoid self-reinforcement,
trim fraction 0.2, and any single recommender's influence capped at
$\kappa = 0.3$. The cap bounds the effect of a colluding recommender, following
the robustness principle of attack-resistant location
estimation~\cite{liu2008attack}: no minority of recommenders can move the
aggregate arbitrarily.

\subsubsection{Historical trust}

$T^H$ is an exponentially weighted average of prior combined trust for the
$(j,i)$ pair with asymmetric decay. The update is applied after the
combination of~\eqref{eq:combined} for window $t$ has been formed, and the
value it contributes to that combination is the one carried forward from the
previous window:
\begin{equation}
\label{eq:th}
T^H_{ij}(t) = \lambda\, T^H_{ij}(t-1) + (1-\lambda)\, T_{ij}(t),
\qquad
\lambda =
\begin{cases}
\lambda_\downarrow = 0.5, & T_{ij}(t) < T^H_{ij}(t-1),\\
\lambda_\uparrow = 0.9, & \text{otherwise,}
\end{cases}
\end{equation}
so trust degrades faster than it recovers: wrongly trusting an attacker costs
more than briefly distrusting an honest node. On the first window for a pair,
$T^H$ is undefined, drops out of the renormalisation, and is seeded from the
resulting combination.

\subsubsection{Combined trust}

\begin{equation}
\label{eq:combined}
T_{ij}(t) = \frac{w_D T^D_{ij}(t) + w_I T^I_{ij}(t) + w_H T^H_{ij}(t-1)}{w_D + w_I + w_H},
\end{equation}
renormalised over whichever terms are defined in the window. Because $T^H$
seeds from the combined value, it cannot be evaluated in isolation: a
configuration weighting history alone produces no defined trust at any window.
This constrains the ablation design of Section~\ref{sec:ablation}.

\subsection{Confidence, state and confirmation}

Per-observer trust is collapsed to a single per-target confidence by the median
over observers holding a defined opinion in that window,
\begin{equation}
\label{eq:median}
C_i(t) = \operatorname{median}\{\, T_{ij}(t) : j \in O_i(t) \,\},
\end{equation}
where $O_i(t)$ is that observer set and $|O_i(t)| \ge 3$ is required for any
verdict. The median prevents a single observer with an extreme value from
moving the collapsed confidence. The confidence may optionally be smoothed,
$\hat C_i(t) = \gamma C_i(t) + (1-\gamma)\hat C_i(t-1)$, though the evaluated
configuration uses $\gamma = 1$, so temporal memory enters the decision through
the persistence requirement and not through smoothing. The verdict map is
\begin{equation}
\label{eq:verdict}
S_i(t) =
\begin{cases}
\textsc{Normal}, & \hat C_i(t) \ge \tau_N,\\
\textsc{Suspicious}, & \tau_M < \hat C_i(t) < \tau_N,\\
\textsc{Confirmed}, & \hat C_i(t) \le \tau_M,
\end{cases}
\end{equation}
with $0 \le \tau_M < \tau_N \le 1$. Confirmation also requires
persistence,
\begin{equation}
\label{eq:persist}
M_i(t) = \sum_{k=0}^{W-1} \mathbb{I}\bigl[\hat C_i(t-k) \le \tau_M\bigr],
\end{equation}
with $W = K = 3$, so beacon $i$ is confirmed only if the three most recent
windows all satisfy the condition. A verdict also requires agreement: at
least a fraction $a_{min} = 0.5$ of observers holding a defined opinion, and no
fewer than three observers, must place the target below $\tau_M$ in the same
window.

Two further rules complete the decision logic.

\textbf{Single-branch override.} An observer whose direct trust falls below
$\tau_X = 0.10$, corresponding to $|Z| \ge 2.146$, can confirm a target without
satisfying the agreement fraction. An attack visible to only one evidence class
produces one branch far into its tail while the other stays at the null, and
requiring majority agreement there would suppress the detection the
architecture is meant to enable. The override is load-bearing for the coverage
result: under A1 the behavioural branch contributes nothing, and confirmation
proceeds on localisation evidence alone. The A1 and A2 results should be
attributed to fusion together with this override.

\textbf{Asymmetric clearing.} Leaving the confirmed state requires $K_c = 12$
consecutive clean windows against the $K = 3$ required to enter it, a four-fold
asymmetry in the same direction as the trust decay constants. This rule gives
policy E its release behaviour.

\textbf{Evidence-aware release.} Isolation removes a beacon from forwarding, so
while it is isolated its behavioural branch receives too few transit packets
to contribute, and its trust rests on localisation alone. Under policy E a
beacon confirmed on behavioural evidence can therefore be cleared by the
absence of that evidence rather than by its return. Policy E$'$ closes this
path: a beacon confirmed in a window in which behavioural evidence was
available is released only on a clean window in which behavioural evidence is
again available. A beacon confirmed on localisation evidence alone is released
as under E. Section~\ref{sec:release} measures what the rule costs.

\subsection{Isolation}

On confirmation the beacon enters the isolation set, and under a releasing
policy it leaves when released,
\begin{equation}
\label{eq:iso}
\mathcal{B}_{iso}(t) = \bigl(\mathcal{B}_{iso}(t-1) \cup \{\, i : S_i(t) = \textsc{Confirmed} \,\}\bigr)
\setminus \mathcal{R}(t),
\end{equation}
where $\mathcal{R}(t)$ is the set released at $t$ (empty for the policies
that never release). The admissible node set becomes $V'(t) = V \setminus \mathcal{B}_{iso}(t)$.
Every subsequent route must satisfy $P \cap \mathcal{B}_{iso}(t) = \emptyset$.
Enforcement is by a blacklist queueing discipline; route repair is performed by
AODV.

\subsection{Evaluation quantities}

Recovery time is $T_{rec} = T_{new\text{-}route} - T_{isolation}$. Standard
network measures are
\begin{equation}
\label{eq:netm}
PDR = \frac{P_{received}}{P_{sent}}, \qquad
TH = \frac{8 D_{received}}{T_{sim}}, \qquad
RO = \frac{N_{control}}{N_{data}},
\end{equation}
and we report the throughput recovery ratio
\begin{equation}
\label{eq:recov}
R_{TH} = \frac{TH_{after}}{TH_{normal}},
\end{equation}
for which a well-behaved mechanism gives $R \to 1$.

Finally we define a false-isolation cost that we believe is under-reported in
this literature. Let $H$ be the set of honest beacons and $\iota_i(t)$ the
indicator that $i$ is isolated at time $t$. Cumulative false isolation over a
run of duration $T$ is
\begin{equation}
\label{eq:Phi}
\Phi = \int_0^T \sum_{i \in H} \iota_i(t)\, dt,
\end{equation}
in node-seconds. $\Phi$ measures service withdrawn from the network by the
security mechanism itself. A mechanism with a low false-positive rate may
still have large $\Phi$ if it never releases. Because $\Phi$ grows with the
number of honest beacons and with run duration, we also report the
dimensionless form
\begin{equation}
\label{eq:Phibar}
\bar\Phi = \frac{\Phi}{|H| \cdot T},
\end{equation}
the mean fraction of honest-beacon service time withdrawn by the mechanism,
which is comparable across topologies.
